\setcounter{subsection}{0}
\section{CONTEXTO DEL PROYECTO}

Los capítulos anteriores plantearon el problema y fijaron los cinco objetivos específicos y la
metodología híbrida. Antes de recoger requerimientos hacía falta conocer el terreno. La compra de
útiles escolares parece una tarea sencilla de comercio, pero automatizarla exige resolver problemas
de campos distintos: la lectura de documentos sin formato común, la interpretación de un vocabulario
que cada colegio escribe a su manera, la elección entre muchas combinaciones de productos y la
presentación de las opciones sin abrumar a quien decide. Este capítulo se apoya en la etapa de
comprensión del negocio de CRISP-DM \citep{chapman_crisp-dm_2000} y en el análisis del entorno de
Design Thinking \citep{brown_design_2008}. Primero presenta los retos de la lista escolar y las
técnicas que atiende cada uno (Tabla~\ref{tab:antecedentes}); luego revisa las soluciones existentes
y las alternativas consideradas, y las compara con la propuesta. El vacío que resulta, entre las
soluciones actuales y una recomendación completa a partir de la lista del colegio, es el punto de
partida del Capítulo~IV.\par

\subsection{Antecedentes}

En esta sección se presentan los conceptos necesarios para comprender el funcionamiento de
Tukit. Se parte de la lista escolar y de los retos que plantea convertirla en una compra, y se
asocia a cada reto la técnica que lo resuelve.\par

\subsubsection{La lista escolar y los retos de convertirla en una compra}

La lista de útiles es el documento con el que un colegio indica a las familias qué
materiales necesita cada estudiante para el año escolar. Aunque su propósito es el mismo en
todas las instituciones, cada colegio la elabora a su manera. Algunas listas se publican
como tablas con varias columnas, otras como textos con viñetas, y muchas llegan a las
familias como fotografías compartidas en grupos de mensajería. Además, el mismo artículo se
nombra de formas distintas, como ``cartulina en octavos de colores'' o ``cartulina octavos
surtida'', mientras que nombres casi iguales, como ``block oficio'' y ``block carta'',
corresponden a productos diferentes.\par

Convertir esa lista en una compra completa implica superar varios retos sucesivos. El
primero es leer el documento sin perder renglones, sin importar su formato ni la calidad de
la imagen. El segundo es entender qué pide cada renglón, a pesar de los errores de
ortografía, de los sinónimos y de la diferencia entre la cantidad que solicita el colegio y
el empaque en que se vende el producto. El tercero es encontrar en el catálogo el producto
que corresponde, sin ofrecer uno parecido pero equivocado, ya que la familia solo descubre
ese error cuando ya ha comprado. El cuarto es elegir, entre todas las combinaciones de
productos posibles, aquellas que mejor equilibran el costo, la calidad y las preferencias de
la familia. Finalmente, el quinto es presentar esas alternativas de manera que faciliten la
decisión en lugar de dificultarla.\par

Cada reto se apoya en un cuerpo de conocimiento. La Tabla~\ref{tab:antecedentes} los resume con
la idea que Tukit toma de cada uno, y el Anexo~\ref{anx:antecedentes} los desarrolla. La apuesta
del proyecto es combinarlos en una cadena de lectura, normalización, optimización y revisión humana.

\begin{table}[htbp]
\centering
\scriptsize
\caption{Conceptos de base, el reto que atienden y lo que Tukit toma de cada uno.}
\label{tab:antecedentes}
\begin{tabularx}{\textwidth}{@{}p{3.2cm}p{2.2cm}X@{}}
\toprule
Concepto & Reto & Lo que toma Tukit \\
\midrule
Sistemas de recomendación basados en conocimiento \citep{adomavicius_toward_2005,burke_hybrid_2002} &
Recomendar & Recomienda desde la lista oficial y las preferencias declaradas, no desde lo que compraron
otras familias; recomienda paquetes completos y coherentes \citep{avny_brosh_bruce_2022} \\
Comprensión de documentos con modelos multimodales \citep{shen_ocr_2026,liu_ocrbench_2024} & Leer &
Lee foto, PDF o texto razonando sobre la disposición de la página, con una salida de formato fijo que
se verifica antes de aceptarse \\
Representación vectorial y búsqueda híbrida \citep{reimers_sentence-bert_2019,karpukhin_dense_2020} &
Entender y encontrar & La búsqueda léxica decide qué conceptos son admisibles y la semántica desempata
o reconoce sinónimos; de RAG \citep{lewis_retrieval-augmented_2020} toma la indexación y la
recuperación, no la redacción \\
Optimización multiobjetivo y NSGA-II \citep{ehrgott_multiobjective_2002,deb_fast_2002} & Elegir & Con
$3^{n}$ combinaciones, aproxima el frente de Pareto y mide su calidad con el hipervolumen
\citep{zitzler_multiobjective_1999}; de él toma un kit por perfil \\
Revisión humana (\emph{human-in-the-loop}) \citep{amershi_guidelines_2019} & Evitar errores & Los
renglones inciertos pasan a una cola del equipo antes de llegar a las familias \\
Carga cognitiva y arquitectura de la decisión \citep{sweller_cognitive_2011,iyengar_when_2000,thaler_nudge_2008}
& Presentar & Pocos kits bien diferenciados, cada uno con su razón, y cualquier producto se puede
cambiar antes de pagar \\
\bottomrule
\end{tabularx}
\end{table}

\subsection{Análisis del Contexto}
\label{sec:analisis-contexto}

\subsubsection{Soluciones existentes}

A nivel internacional, varias plataformas han buscado simplificar la compra de útiles
escolares a partir de la lista que publica cada institución. En Estados Unidos, TeacherLists
reúne las listas de cerca de dos millones de salones de clase y las pone a disposición de
cadenas minoristas, que pueden convertir cada lista en un carrito de compra prellenado
\citep{teacherlists_retail_2026}. Por su parte, 1st Day
School Supplies vende kits de útiles a las familias de los colegios aliados, con fechas límite
de pedido y con marcas reconocidas \citep{firstday_faqs_2026}. Estas soluciones confirman
que pasar de la lista al carrito de forma automática es viable y que las familias valoran
esa comodidad. Sin embargo, todas dependen de que la lista se registre manualmente en su
plataforma con un formato estructurado, ofrecen una sola forma de completarla y no están
disponibles en Colombia. En el caso de 1st Day School Supplies, además, solo pueden
comprar las familias de colegios aliados, dentro de fechas límite de pedido y con las
marcas que ofrece la empresa.\par

En el contexto colombiano, las alternativas disponibles tampoco parten de la lista oficial
de cada colegio. Panamericana y Éxito organizan su oferta escolar por categorías o por
nivel educativo \citep{panamericana_cotizador_2025,exito_listas_2025}, Primavera permite
que el acudiente arme manualmente su lista sobre el catálogo propio
\citep{primavera_armatulista_2026}, y MercadoLibre comercializa kits fijos que arma cada
vendedor \citep{mercadolibre_kits_2025}. En todos los casos, la familia debe traducir por su
cuenta cada renglón de la lista en un producto, comparar alternativas y verificar que la
compra quede completa, lo que mantiene la carga cognitiva del proceso tradicional
\citep{sweller_cognitive_2011}.\par

En el ámbito académico, la recomendación de paquetes estudia cómo sugerir conjuntos de
productos coherentes entre sí, por ejemplo mediante representaciones vectoriales de los
artículos que componen cada paquete \citep{avny_brosh_bruce_2022,du_enhancing_2025} o
mediante modelos de lenguaje que generan o evalúan paquetes completos
\citep{liu_llmcbr_2025,bayrak_evaluating_2024}. De forma paralela, la generación aumentada
por recuperación se ha aplicado al comercio electrónico para responder preguntas sobre
productos \citep{xu_itemrag_2026}. Estos trabajos aportan técnicas valiosas, pero suponen un
historial de interacciones o una consulta escrita por el usuario, y no abordan el caso en
el que el punto de partida es un documento institucional sin formato fijo que debe
cumplirse completo.\par

\subsubsection{Alternativas de solución consideradas}

Ante este panorama, se analizaron seis formas de resolver el problema, según su viabilidad técnica
y su efecto sobre las familias y la operación. Las tres primeras abordan el problema completo; las
tres últimas resuelven solo la lectura de la lista.\par

\begin{enumerate}
  \item \textbf{Listas digitalizadas manualmente.} Siguiendo el modelo de TeacherLists, el
        colegio o el equipo de Tukit transcribe cada lista en un formulario estructurado, y
        cada renglón se asocia de antemano con un producto.
        \begin{itemize}
          \item \textit{Fortaleza:} la correspondencia entre renglón y producto es exacta,
                porque la define una persona.
          \item \textit{Debilidad:} exige transcribir cada lista y cada cambio de lista, no
                permite que el acudiente cargue la suya y no escala a cientos de colegios.
        \end{itemize}

  \item \textbf{Kits preconfigurados por colegio y grado.} Siguiendo el modelo de 1st Day
        School Supplies, se arma un único kit fijo para cada lista.
        \begin{itemize}
          \item \textit{Fortaleza:} la compra se resuelve en un solo paso y la operación es
                sencilla.
          \item \textit{Debilidad:} ofrece una sola alternativa, sin considerar el
                presupuesto ni las preferencias de cada familia, y requiere armar y
                actualizar manualmente cada kit.
        \end{itemize}

  \item \textbf{Generación directa con un modelo de lenguaje.} Un modelo de lenguaje recibe
        la lista y redacta el kit completo, siguiendo el esquema de generación aumentada
        por recuperación \citep{lewis_retrieval-augmented_2020}.
        \begin{itemize}
          \item \textit{Fortaleza:} se adapta con facilidad a listas redactadas de formas
                muy distintas.
          \item \textit{Debilidad:} puede proponer productos, cantidades o precios que no
                están en el catálogo, sus resultados son difíciles de auditar y no garantiza
                el mejor equilibrio entre costo, calidad y preferencias.
        \end{itemize}

  \item \textbf{Reconocimiento óptico de caracteres con reglas y búsqueda exacta.} La lista
        se lee con OCR tradicional \citep{smith_overview_2007} y cada renglón se busca en el
        catálogo por coincidencia exacta de palabras.
        \begin{itemize}
          \item \textit{Fortaleza:} es económica y sus resultados son fáciles de explicar.
          \item \textit{Debilidad:} falla con fotografías inclinadas, con tablas de varias
                columnas y con las múltiples formas de nombrar un mismo artículo.
        \end{itemize}

\item \textbf{Análisis de maquetación con Docling.} Docling, desarrollado por IBM,
      combina el reconocimiento de caracteres con modelos que identifican la estructura
      del documento, como títulos, párrafos, tablas y columnas, y la convierte en un
      formato organizado \citep{auer_docling_2024}.
      \begin{itemize}
        \item \textit{Fortaleza:} recupera con fidelidad las tablas y las columnas de los
              documentos PDF con una estructura definida, es de código abierto y puede
              ejecutarse en la infraestructura propia.
        \item \textit{Debilidad:} su desempeño disminuye con las fotografías tomadas con
              el celular, que son frecuentes entre las listas que comparten las familias
              (Sección~\ref{sec:percepcion} compara los motores de lectura), y no interpreta el contenido. Por ejemplo, no reconoce que ``Cartuchera con
              esferos y borrador'' describe tres artículos que se compran por separado, de
              modo que la separación de renglones y su normalización quedarían por
              resolver.
      \end{itemize}
        \item \textbf{Plataformas comerciales de análisis documental.} Servicios como Amazon
      Textract \citep{aws_textract_2026} e IBM Watson Discovery \citep{ibm_discovery_2026}
      extraen el texto, las tablas y los campos de un documento en la nube, y permiten
      buscar información dentro de grandes colecciones de archivos.
      \begin{itemize}
        \item \textit{Fortaleza:} ofrecen alta precisión en la extracción de texto y
              tablas, y no requieren desarrollar ni mantener un modelo propio.
        \item \textit{Debilidad:} su costo crece con cada página procesada y los documentos
              se tratan en la infraestructura de un tercero. Además, solo resuelven la
              lectura, ya que no interpretan el vocabulario escolar ni asocian cada
              renglón con un producto del catálogo, de modo que la normalización y la
              construcción de los kits quedarían por resolver.
      \end{itemize}
\end{enumerate}



Frente a estas alternativas, la solución del Capítulo~II toma de cada una lo que resuelve una
parte del problema. Conserva la lista oficial del colegio como punto de partida, igual que
las listas digitalizadas, pero la lee de forma automática en cualquier formato mediante un
modelo multimodal \citep{shen_ocr_2026}. Aprovecha la flexibilidad de los modelos de
lenguaje, pero solo para leer la lista, mientras que la asociación con el catálogo se
resuelve con una búsqueda híbrida que no inventa productos \citep{karpukhin_dense_2020}. En
lugar de un único kit fijo, construye varias alternativas que equilibran el costo, la
calidad y las preferencias mediante optimización multiobjetivo \citep{deb_fast_2002}, y
mantiene la revisión humana en los casos dudosos, lo que permite auditar cada resultado.\par

\subsubsection{Comparación con la solución propuesta}

La Tabla~\ref{tab:comparacion-soluciones} (Anexo~\ref{anx:tablas}) contrasta, en nueve capacidades,
a TeacherLists, 1st Day Supplies, Panamericana, Grupo Éxito, Papelería Primavera y MercadoLibre con
Tukit. Hasta donde se revisó, ninguna de las seis lee listas en foto, PDF o texto ni reconoce el
vocabulario escolar local, y solo Tukit parte de la lista oficial del colegio, combina productos de
varias tiendas y ofrece alternativas según presupuesto y preferencias. Las soluciones comerciales
superan a Tukit en catálogo general y control del inventario, porque Tukit no vende ni almacena. La
brecha que Tukit busca cubrir es, entonces, convertir cualquier lista escolar en kits completos
y auditables, que es lo que ninguna de las soluciones revisadas hace.\par

% Tabla tab:comparacion-soluciones en el Anexo A.3 (tablas del cuerpo)
